\documentclass{article}
\usepackage{xcolor}
\usepackage{titleps}
\usepackage[letterpaper, margin=0.95in]{geometry}
\usepackage{url}
\usepackage{amsmath}
\usepackage{amssymb}
\usepackage{wrapfig}
\usepackage{float}
\usepackage{mathtools}
\usepackage{enumitem}
\usepackage{tabu}
\usepackage{parskip}
\usepackage{natbib}
\usepackage{graphicx}
\usepackage{listings}

\usepackage[many]{tcolorbox}
\usepackage{minted}
\setminted[python]{
	% frame=single,
	% linenos,
    xleftmargin=0.475em,
    baselinestretch=1.2,
}
% https://tex.stackexchange.com/a/569249
\setcounter{secnumdepth}{5}
\setcounter{tocdepth}{5}
\makeatletter
\newcommand\subsubsubsection{\@startsection{paragraph}{4}{\z@}{-2.5ex\@plus -1ex \@minus -.25ex}{1.25ex \@plus .25ex}{\normalfont\normalsize\bfseries}}
\newcommand\subsubsubsubsection{\@startsection{subparagraph}{5}{\z@}{-2.5ex\@plus -1ex \@minus -.25ex}{1.25ex \@plus .25ex}{\normalfont\normalsize\bfseries}}
\makeatother

\usepackage{hyperref}
\usepackage[color=red]{todonotes}
\usepackage{forest}
\definecolor{light-yellow}{HTML}{FFE5CC}

\newpagestyle{ruled}
{\sethead{CMU 16-831}{Introduction to Robot Learning }{Fall 2026}\headrule
  \setfoot{}{}{}}
\pagestyle{ruled}

\renewcommand\makeheadrule{\color{black}\rule[-.75\baselineskip]{\linewidth}{0.4pt}}
\renewcommand*\footnoterule{}

\newtcolorbox[]{answer}[1][]{
    % breakable,
    enhanced,
    nobeforeafter,
    colback=white,
    title=Your Answer,
    sidebyside align=top,
    box align=top,
    #1
}



\begin{document}

\lstset{basicstyle = \ttfamily,columns=fullflexible,
backgroundcolor = \color{light-yellow}
}

\begin{centering}
    {\Large Assignment 3: Q-Learning and Actor-Critic Algorithms
} \\
    \vspace{.25cm}
\end{centering}
\vspace{0.25cm}

\textbf{Andrew ID:} \texttt{jiayuy2} \\
\textbf{Collaborators:} \texttt{None}\\ 
\textbf{NOTE:} Please do \textbf{NOT} change the sizes of the answer blocks or plots.

\setcounter{section}{0}
\section{Question 1: DQN and DDQN (LunarLander-v3) - [7.5 points]}
\begin{answer}[title=Question 1,height=9.5cm,width=\linewidth]
\centering
\includegraphics[height=7cm]{figures/q1.png}

{\footnotesize LunarLander-v3, 3 seeds each, default hyperparameters. Solid: mean \texttt{Train\_AverageReturn} ($\pm$1 std shaded); dashed: mean \texttt{Train\_BestReturn}. DDQN picks $a'$ with the online net and evaluates it with the target net, reducing the max-induced overestimation of DQN; DQN peaks then dips, DDQN is steadier and ends higher.}
\end{answer}

\section{Question 2: Actor-Critic (CartPole-v0) - [3.75 points]}
\begin{answer}[title=Question 2,height=9.5cm,width=\linewidth]
\centering
\includegraphics[height=7.2cm]{figures/q2.png}

{\footnotesize CartPole-v0, \texttt{q2\_10\_10} (\texttt{-ntu 10 -ngsptu 10}, \texttt{-n 100 -b 1000}). \texttt{Eval\_AverageReturn} reaches 200, matching the hw2 policy-gradient result.}
\end{answer}



\section{Question 3: Actor-Critic (InvertedPendulum-v4) - [3.75 points]}
\begin{answer}[title=Question 3,height=9.5cm,width=\linewidth]
\centering
\includegraphics[height=7.2cm]{figures/q3.png}

{\footnotesize InvertedPendulum-v4, \texttt{q3\_10\_10} (\texttt{--ep\_len 1000 --discount 0.95 -n 100 -l 2 -s 64 -b 5000 -lr 0.01 -ntu 10 -ngsptu 10}). \texttt{Eval\_AverageReturn} rises past 100 within $\sim$20 iterations and fluctuates around 1000 by iteration 100.}
\end{answer}







\end{document}